% C1 Inicio del Preámbulo
% C1.1 Definición de la Clase de Documento
\documentclass[12pt, doc]{apa7}

% C1.2 Configuración de Idioma y Fuentes
\usepackage[spanish]{babel}
\usepackage{fontspec}
\setmainfont{Times New Roman}

% C1.3 Geometría y Maquetación de la Página
\usepackage[left=1in,right=1in,top=1in,bottom=1in]{geometry}
\usepackage{setspace}
\parindent=0.5in

% C1.4 Elementos Visuales y Manejo de Color
\usepackage{graphicx}
\usepackage{xcolor}

% C1.5 Hipervínculos y Navegación del Documento
\usepackage{csquotes}
\usepackage{hyperref}
\hypersetup{
	pdfborder={0 0 0},
	colorlinks=true,
	linkcolor=black,
	urlcolor=blue,
	citecolor=black
}

% C1.6 Estructura de Tablas Avanzadas
\usepackage{array}
\usepackage{tabularx}
\usepackage{ltablex}
\keepXColumns
\usepackage{booktabs}

% C1.7 Metadatos y Afiliación Académica
\title{}
\author{}
\shorttitle{}
\affiliation{
	SENA - Centro de Electricidad, Electrónica y Telecomunicaciones \\
	Tgo. en Gestión e Implementación de Bases de Datos \\
	\today
}

% Configuración de bibliografía
\usepackage[style=apa, backend=biber]{biblatex}
\addbibresource{referencias.bib}

% ---------------------- FIN DEL PREAMBULO --------------------------------------------------

\begin{document}
	
	% --- PORTADA MANUAL ---
	\begin{titlepage}
		\centering
		\setstretch{1.5}
		
		{\Large\bfseries GA3-220501106-AA2-EV03. Bitácora en herramientas de administración y supervisión \par}
		
		\vspace{5.5cm}
		
		{\large \textbf{Integrantes:} \par}
		\vspace{0.3cm}
		Angie Paola Sepúlveda Moreno \\
		Jhon Octavio Parra Martínez \\
		John Jairo Duarte Aponte \\
		Jorge Andrés Castellanos Soler \\
		José Roberto Pérez Cano \\
		
		\vfill
		
		SENA - Centro de Electricidad, Electrónica y Telecomunicaciones \\
		Tgo. en Gestión e Implementación de Bases de Datos - 3466363 \\
		
		\vspace{1cm}
		
		\textbf{Docente:} \\
		Ing. Andrés Quevedo \\
		
		\vspace{1.5cm}
		
		25 de agosto de 2026 \par
		
	\end{titlepage}
	% --- FIN DE LA PORTADA ---

	\newpage
	\setstretch{1.15}
	\tableofcontents
	\newpage
	\doublespacing
	
	% --- INTRODUCCIÓN ---
	\section{Introducción}
	El presente documento constituye la bitácora de administración y supervisión de los servicios de Tecnologías de la Información (TI) para la empresa \textbf{Alliance Technology Group}, un BPO de Contact Center localizado en la ciudad de Cali, Valle del Cauca de alta exigencia operativa. En este tipo de organizaciones, la disponibilidad de los servicios, la ciberseguridad y la gestión eficiente de los incidentes son pilares fundamentales para garantizar la continuidad del negocio y el cumplimiento de los Acuerdos de Nivel de Servicio (SLA) establecidos con los diferentes clientes que contratan las campañas \parencite{sena2026redes}.
	
	Para lograr un control óptimo y mitigar los impactos negativos frente a cualquier eventualidad, el departamento de TI emplea un ecosistema robusto de herramientas tecnológicas (Fortigate, ServiceDesk Plus, Snipe-IT, Wazuh, Directorio Activo, Action1, UpTimeRobot y PRTG), las cuales permiten centralizar desde la monitorización del hardware y la red hasta la gestión de identidades y la mesa de ayuda \parencite{sena2026servidores}.
	
	\section{Identificación de Buenas Prácticas de Gestión de Servicios TI}
	Alineados con estándares internacionales como ITIL y los principios de ciberseguridad corporativa, en Alliance Technology Group se aplican las siguientes buenas prácticas en la gestión de servicios TI \parencite{sena2026servidores, sena2026redes}:
	
	\begin{itemize}
		\item \textbf{Gestión de Incidentes y Mesa de Servicio (Service Desk):} Centralización de los reportes y peticiones de los usuarios a través de ServiceDesk Plus para ofrecer tiempos de respuesta medibles, priorizando la reducción del MTTR (Tiempo Medio de Reparación) y garantizando el cumplimiento de los SLA.
		\item \textbf{Gestión Proactiva y Monitoreo de la Infraestructura:} Uso de herramientas de monitorización continua como Paessler PRTG y UpTimeRobot para identificar latencias, caídas de servicios y degradación de red antes de que afecten la operación de los agentes en el Contact Center.
		\item \textbf{Gestión de Activos de TI (ITAM):} Trazabilidad y control del ciclo de vida del hardware (diademas, estaciones de trabajo) y licencias de software mediante Snipe-IT.
		\item \textbf{Gestión de Ciberseguridad (SGSI):} Auditoría continua de bitácoras y logs, detección de vulnerabilidades e intrusiones mediante Wazuh (SIEM) y protección perimetral estricta con appliances Fortigate.
		\item \textbf{Gestión de Parches y Mantenimiento Predictivo:} Automatización del despliegue de parches de seguridad y actualizaciones de software mediante Action1 para mitigar riesgos cibernéticos y operacionales.
		\item \textbf{Control de Acceso y Gestión de Identidades (IAM):} Administración centralizada de credenciales de usuario, autenticación de red y políticas de seguridad (GPOs) mediante el Directorio Activo de Microsoft Windows Server.
	\end{itemize}
	
	\section{Bitácora de Herramientas de Administración y Supervisión}
	A continuación, se detalla la bitácora técnica de supervisión, en cumplimiento de los criterios de desempeño sobre infraestructura TI, herramientas de administración y gestión bajo demanda. La información a continuación es suministrada al equipo de trabajo por \textbf{José Roberto Pérez Cano} bajo el rol de Analista de Soporte de Bases de Datos y Aplicaciones adscrito al área de IT de la empresa.
	
	\vspace{0.3cm}
	\noindent\textbf{Datos Generales de la Bitácora:}
	\begin{itemize}
		\item \textbf{Nombre de la organización:} Alliance Technology Group (Contact Center BPO)
		\item \textbf{Nombres de los aprendices:} Angie Paola Sepúlveda Moreno, Jhon Octavio Parra Martínez, John Jairo Duarte Aponte, Jorge Andrés Castellanos Soler, José Roberto Pérez Cano
		\item \textbf{Fecha de registro:} 25 de agosto de 2026
		\item \textbf{Área o dependencia de la organización:} Área de Sistemas o IT Department (Nombre Oficial Institucional)
	\end{itemize}
	
	\vspace{0.5cm}
	\noindent
	\begin{tabularx}{\textwidth}{|p{3.5cm}|X|X|}
		\hline
		\textbf{Herramienta} & \textbf{Descripción Técnica (Servicio / Plataforma)} & \textbf{Funciones de Administración y Supervisión Asignadas} \\ \hline
		\textbf{Fortigate} & Firewall de Próxima Generación (NGFW) y gestión unificada de amenazas perimetrales. & Supervisión y control del ancho de banda, configuración de VPNs para agentes en teletrabajo, bloqueo de intrusiones (IPS), filtrado de contenido web y detección de malware de borde. \\ \hline
		\textbf{Sophos} & Solución de seguridad endpoint y red. & Gestión de antivirus, detección y respuesta para puntos finales, y protección avanzada contra ransomware e intrusiones. \\ \hline
		\textbf{ServiceDesk Plus} & Software especializado de Mesa de Ayuda (Help Desk) y Gestión de Servicios (ITSM). & Recepción, categorización y seguimiento de tickets de soporte técnico, gestión del catálogo de servicios, cumplimiento de SLA y métricas de rendimiento por agente de soporte. \\ \hline
		\textbf{Snipe-IT} & Sistema Open Source para la gestión y seguimiento de activos de TI (ITAM). & Administración rigurosa del inventario físico (PCs, monitores, periféricos) y lógico (licencias), control de asignaciones, mantenimientos realizados y seguimiento de garantías. \\ \hline
		\textbf{Wazuh} & Plataforma SIEM (Security Information and Event Management) y XDR. & Recolección centralizada de logs, monitoreo de integridad de archivos críticos, detección de rootkits, escaneo de vulnerabilidades y alertas automatizadas ante posibles brechas de seguridad. \\ \hline
		\textbf{Directorio Activo (Active Directory)} & Servicio de directorio de Microsoft Windows Server para gestión de redes de dominio. & Gestión centralizada del ciclo de vida de identidades, autenticación, agrupación de perfiles por áreas y despliegue masivo de políticas de seguridad restrictivas (GPO) en estaciones. \\ \hline
		\textbf{Action1} & Plataforma basada en la nube para la gestión de endpoints y despliegue de parches (RMM). & Inventariado de software, despliegue forzado de actualizaciones del SO y terceros, ejecución remota de scripts de corrección y soporte remoto de bajo nivel a los equipos. \\ \hline
		\textbf{UpTimeRobot} & Herramienta SaaS (Software as a Service) para el monitoreo del tiempo de actividad web. & Supervisión continua cada minuto de portales críticos (CRM, pasarelas SIP), APIs y sitios externos, emitiendo alertas inmediatas por correo/SMS ante caídas para una rápida contingencia. \\ \hline
		\textbf{Paessler PRTG} & Software integral de monitorización y sensores de infraestructura de red unificada. & Supervisión del estado de salud de los servidores (uso de CPU, RAM, espacio en discos), sensores de latencia, monitorización de tráfico en switches y disponibilidad de los recursos LAN/WAN. \\ \hline
		\textbf{Proxmox} & Entorno de virtualización de código abierto. & Administración de máquinas virtuales (VMs) y contenedores Linux (LXC), gestión de clústeres, almacenamiento y copias de seguridad. \\ \hline
		\textbf{Docker} & Plataforma de contenedores. & Creación, despliegue y ejecución de aplicaciones en contenedores aislados, optimizando la portabilidad y la eficiencia de los recursos. \\ \hline
		\textbf{Portainer} & Herramienta de gestión de contenedores con interfaz gráfica. & Supervisión y administración centralizada de contenedores Docker, imágenes y redes, simplificando el despliegue y mantenimiento de los servicios virtualizados. \\ \hline
		\textbf{Git / GitHub} & Sistema de control de versiones distribuido y plataforma de alojamiento. & Seguimiento de cambios en el código y archivos de configuración, colaboración en equipo y automatización de despliegues (CI/CD). \\ \hline
		\textbf{NGINX Proxy Manager} & Interfaz de administración para proxies inversos basados en NGINX. & Gestión del enrutamiento de tráfico web hacia diferentes servicios internos, configuración de certificados SSL/TLS y control de acceso. \\ \hline
	\end{tabularx}
	\vspace{0.5cm}
	
	\section{Conclusiones}
	La consolidación y ejecución de esta bitácora demuestra que Alliance Technology Group posee un esquema maduro de administración TI, abarcando eficientemente desde la mesa de servicio hasta la ciberseguridad perimetral y la gestión de infraestructura moderna. La utilización conjunta e integrada de estas herramientas permite a la organización reaccionar rápidamente ante incidentes, aplicar un enfoque estructurado de carácter preventivo y asegurar una alta tolerancia a fallos, confiabilidad técnica y confidencialidad total sobre la información de las diferentes campañas \parencite{sena2026servidores}.
	
	Además, la implementación de estas buenas prácticas y herramientas permite concluir lo siguiente \parencite{sena2026servidores, sena2026redes}:
	
	\begin{itemize}
		\item \textbf{Mejora en los Niveles de Servicio:} El monitoreo constante con herramientas como PRTG y UpTimeRobot facilita el cálculo preciso de indicadores como el MTTR (Tiempo Medio de Reparación) y el MTBF (Tiempo Medio Entre Fallas). Un MTTR bajo garantiza que el equipo de TI responda rápidamente a los incidentes, mientras que un MTBF alto refleja la fiabilidad de los activos y la eficacia de los mantenimientos.
		\item \textbf{Mantenimiento Predictivo Eficaz:} La supervisión y análisis de datos permiten implementar planes de mantenimiento preventivo y prescriptivo. Esto disminuye los costos operativos, prolonga la vida útil de los equipos y reduce significativamente los tiempos de inactividad imprevistos, minimizando el impacto en la operación del Contact Center.
		\item \textbf{Agilidad y Eficiencia con Contenedores:} El uso de plataformas como Docker y Portainer aporta agilidad, seguridad y menor sobrecarga al sistema. Al virtualizar a nivel del sistema operativo, se facilita el desarrollo, despliegue y actualización de aplicaciones, optimizando los recursos de hardware y garantizando una mayor eficiencia en los proyectos de TI.
		\item \textbf{Beneficios del Cloud Computing:} La adopción de soluciones en la nube, gestionadas a través de la infraestructura descrita, brinda ventajas competitivas como alta disponibilidad, resiliencia ante desastres, seguridad reforzada de la información, y la capacidad de acceso remoto, factores críticos para el teletrabajo y la continuidad del negocio.
	\end{itemize}
	
	\nocite{*} 
	\printbibliography[title={Referencias}]
	
\end{document}